\part{Questions the professor may ask}

\section{Questions and answers}
\label{sec:qa}

Short answers we can say out loud, with the section that has the details.

\subsection{Motivation and novelty}

\Q{What is new here compared with SQLancer?}{SQLancer's oracles (TLP, NoREC,
PQS) run on queries from its own small generators, whose structure it knows.
We apply the same and more oracles to \emph{derivation trees of a large typed
grammar} (78 rules, 421 alternatives), at any level of a nested query, and we
use the same tree for mutation and for reduction. Two further things SQLancer
does not do: steering towards recent commits with a gated LLM, and measuring
oracle precision with a ledger. And a ground-truth benchmark of 23 re-injected
\pg bugs to compare generators fairly.}

\Q{Why not just use SQLsmith or SQLancer as the baseline?}{The paper compares
\emph{generators} under identical oracles and budget; plugging an external tool
into our executor would change two things at once. Our random/mutation arm is
the standard unguided baseline (AFL-style text mutation). We discuss SQLancer,
SQLsmith and SQLRight in related work, and the motivating example (BUG-02) is
a bug that had not been reported although it has been in released \pg since 17.}

\Q{Why \pg and why master?}{\pg is the most used open-source DBMS; its master
branch gets new planner optimisations every month, and regressions there are
fixed fastest (``open items'' before a release), which helps our chance of an
accepted bug. We also check every finding on the 18.6 release: BUG-01 and
BUG-02 affect released versions.}

\Q{Why do wrong results matter more than crashes?}{A crash is loud: the
user sees an error and the log has a trace. A wrong result is silent: an
application gets wrong rows and nobody notices. 11 of the 20 observable
benchmark bugs are invisible to a crash-only fuzzer (Section~\ref{sec:bench}).}

\Q{In one sentence, what is ``one tree, four jobs''?}{The same typed
derivation tree generates the SQL, is mutated (regenerate, splice, insert,
drop subtrees), is rewritten by the oracles into an equivalent query, and is
shrunk when a test case fails; every job reuses the grammar's knowledge of
types and scopes.}

\subsection{Grammar and generator}

\Q{How did validity go from 83\% to 99.5--100\%?}{By a measurement loop:
\texttt{pgfuzz sample --execute} prints the valid rate, top server error
messages and the alternatives with the worst valid rate; each fix targets one
rule (type a function's arguments correctly, keep a column in scope, avoid a
million-row \texttt{generate\_series}). Four rounds took it from 83.4\% to
100\% on 966 statements (Section~\ref{sec:story}). With the larger grammar
today it is 99.5--99.7\%.}

\Q{Why is the grammar typed? Isn't a CFG enough?}{A CFG gives syntactically
valid SQL that is mostly rejected by the analyzer: \texttt{int + text}, unknown
columns, aggregates in \texttt{WHERE}. Types and scopes make generation a
proof search for $\Gamma \vdash e : T$ (L29--30), so 99.5\% of statements
reach the planner and executor, which is where the bugs are. Types are also
what makes rewrites sound: \texttt{e+0} is an identity for \texttt{int}, not
for \texttt{float8} ($-0 + 0$ prints 0).}

\Q{What is the derivation tree exactly?}{A node per rule expansion, recording
which alternative was chosen, the bound type, and for builtins the chosen
value (a column, a literal). Its leaves read left to right are the SQL. The
renderer wraps every node's tokens in invisible markers, so we know which
token span each node produced.}

\Q{How does replay work, and why is it needed?}{Given a stored tree, the
engine re-derives it following the same choices, but re-binds names and types
in the current state. Mutation and reduction edit a tree and replay it, so the
result stays valid even if aliases shift. Trees are stored as JSON with stable
alternative ordinals; 599 of 600 test cases replay to identical SQL in a fresh
process.}

\Q{What does the bandit do and is it significant?}{Multiplicative weights
over grammar alternatives: reward 1 if the test case entered the corpus (new
signature, feature, feature pair or plan shape), $-0.5$ if the statement
failed; weight $\times e^{\eta(r-\bar r)}$, $\eta = 0.1$, clipped to $[0.2,5]$.
In one 1{,}000-test-case run: 533 vs 419 plan shapes at the same validity. One
run is not significance; RQ1 has an ablation arm (\texttt{grammar\_mm\_flat})
with 5 trials.}

\Q{What is the data profile for?}{Constants chosen blindly make predicates
select all rows or none, so the planner never faces a hard choice. The
profile gives per-column min/max/quantiles; \texttt{@dlit} draws a literal from
the column being compared. Predicates selecting 5--95\% of rows went from
29.3\% to 58.6\%. Boundary literals (min$\pm$1, $-0$, NaN, $\pm$Infinity, empty
array) are interval analysis applied to test inputs (L19--20).}

\Q{How big is the schema and why so small?}{Nine relations, about 200 rows
each, deterministic data with NULLs, all index kinds, partitioning,
inheritance, a view, generated columns. Small tables keep each query fast (17--23
test cases per second with all oracles), and most optimizer bugs need
structure (an index, a partition, a NULL), not volume. Volume-dependent bugs are
an external threat.}

\subsection{Oracles}

\Q{How do you know a metamorphic mismatch is a real bug and not your
rewrite?}{Three layers. (1) Each rewrite has a soundness argument and is
restricted by the tree to where it holds (Section~\ref{sec:mm}): multiset
contexts only, unique-text types for \texttt{DISTINCT}, no outer joins for the
NOT NULL conjunct, etc. (2) A mismatch is confirmed by re-running both sides.
(3) Triage reduces it and a human decides; every verdict goes into the
ledger. Where we were wrong (TLP-HAVING, numeric scale), we fixed the class
and added a test, and we report the precision before and after.}

\Q{Explain TLP in one breath.}{For any predicate $p$, every row makes $p$
TRUE, FALSE or NULL, exactly one of them, so the rows of a query equal the
union of the rows with \texttt{WHERE p}, \texttt{WHERE NOT p} and
\texttt{WHERE p IS NULL}. The optimizer treats the four queries very
differently, so a bug in one shows as a mismatch.}

\Q{Explain NoREC.}{\texttt{SELECT count(*) FROM f WHERE p} lets the optimizer
use $p$ (indexes, pushdown, join pruning).
\texttt{SELECT sum(CASE WHEN p THEN 1 ELSE 0 END) FROM f} forces it to evaluate
$p$ on every row. Both count the rows where $p$ is TRUE.}

\Q{Why did TLP-HAVING give false positives?}{Without \texttt{GROUP BY}, a
\texttt{HAVING} clause makes the level a single-group aggregate: it returns one
row even if \texttt{WHERE} removes every row. Our left side \emph{removed}
\texttt{HAVING}, turning it into a plain query: 0 rows vs 1. Now the left side
uses \texttt{HAVING TRUE}. It is a good example of a rewrite that is correct
for the textbook case and wrong in a corner the larger grammar reaches.}

\Q{Why do you mask floats? Doesn't that hide bugs?}{Summing floats in a
different order (hash vs sort aggregation, parallel workers) legitimately
changes the last digits, so exact comparison would flood us with false
positives. We mask only when the query aggregates or uses windows, and the mask
keeps the class (NaN, $\pm$Infinity, finite): when masking everything we lost a
real benchmark bug (0 vs NaN), which taught us this.}

\Q{What does the plan-variant oracle catch that metamorphic rewrites
do not?}{Bugs in one planner path or executor node: a join method, hash
aggregation spilling with tiny \texttt{work\_mem}, parallel plans, partition-wise
join, memoize, eager aggregation. The rewrite oracles catch bugs in transformations that
are always applied (no setting turns them off), like BUG-02. Localization uses
both together.}

\Q{What does the cross-version oracle add?}{An optimisation new on master that
is wrong under every planner setting is invisible to plan variants but
disagrees with 18.6, which does not have it. It needs exclusions for clock
functions and intended changes (\texttt{corr} on master).}

\Q{What is the RLS oracle and why is it ``security''?}{The query runs as a
restricted role with row-level security policies, and as superuser on
explicitly filtered tables; both must be equal. More rows on the restricted
side is a confidentiality leak: the policy is a label and the check is
non-interference on concrete data (L21--22). RLS bugs have had CVEs in \pg
(e.g.\ CVE-2023-2455, CVE-2024-10976). 532 checks, 0 false positives so far.}

\Q{What is CERT?}{Cardinality Estimation Restriction Testing (Ba \& Rigger,
ICSE 2024): adding a \texttt{WHERE} must not raise the planner's row estimate.
A violation is a \emph{performance} outcome, reported separately, never as a
wrong result.}

\Q{Which queries are excluded from comparison and why?}{Those with no single
correct multiset: \texttt{LIMIT/OFFSET} without full ordering, volatile functions,
order-dependent windows, order-sensitive aggregates without \texttt{ORDER BY},
\texttt{ARRAY(SELECT)} without \texttt{ORDER BY}, \texttt{DISTINCT ON}, row
locking, DML. The grammar writes order-sensitive aggregates with their own
\texttt{ORDER BY}, so most queries stay eligible.}

\Q{How many metamorphic checks run per second?}{In the full-system smoke,
1{,}991 checks in 180 s on one worker (about 11/s), alongside 7{,}956
plan-variant runs.}

\subsection{Triage and bugs}

\Q{Walk us through BUG-02.}{Section~\ref{sec:bug02}. Short: since \pg 17 a
planner support function simplifies \texttt{range @> elem} with a constant
range; for an empty or unbounded range it returns the constant false/true and
forgets that a NULL element must give NULL. Our expression rewrite wrapped the
range in \texttt{(SELECT ...)}, which blocks the simplification, and the
results differed for NULL-extended rows of a left join. Fix: only simplify when
the element is provably non-null.}

\Q{Walk us through BUG-01.}{Section~\ref{sec:bug01}. JSON\_TABLE's context
item and PASSING arguments never get collations assigned, so a row comparison
inside them compiles zero column comparisons: false on 18.6, an assertion on
master; text comparisons error. One-line fix: \texttt{assign\_expr\_collations}.}

\Q{Have you reported them? Why not yet?}{Both have reproducers, root causes,
patches and regression tests in \file{docs/bughunt/}. We are filing them this
week (Phase 4 of the plan) after re-running on the newest 17.x/18.x/master
and searching the archives for duplicates (none found on 23 Oct). We wanted
the patches tested with the full regression suite first, which is done.}

\Q{Is BUG-02 really a bug, or intended?}{The operators are declared strict, and
without the planner simplification (the sub-select form) \pg itself returns
NULL. The same expression giving different answers depending on whether a
sub-select surrounds a constant cannot be intended.}

\Q{Why reduce on the tree if you have delta debugging?}{Token reduction gets
stuck: removing any single token breaks the SQL, while a smaller valid query
needs a \emph{different} derivation. Tree reduction replaces a subtree by the
smallest derivation of the same rule and type, and replay keeps it valid. On
the BUG-01 case: 79 tokens vs 116 for token reduction alone, at a similar
number of server runs (555 vs 406).}

\Q{How do you reduce a metamorphic pair without breaking the
equivalence?}{The two queries share a prefix and suffix; we reduce the shared
skeleton with a hole where they differ, and fill the hole with each side's
middle. Every edit applies to both sides at once.}

\Q{How do you deduplicate?}{Signatures, strongest first: top-5 stack frames,
assertion file and expression, SQLSTATE + normalised message, oracle/variant +
statement features, and a weak signal-only key. The strength is logged
because counting by signature is a threat to validity.}

\Q{What is your oracle precision?}{18 triaged candidates so far: 3 real
(BUG-01, BUG-02, one re-injected bug), 13 false positives or intended changes,
2 flaky. Raw precision 0.19, but all 13 came from 6 classes we have since
fixed, and since then 3{,}000+ new-kind, RLS, CERT and full-system checks gave 0
false positives. The long campaigns will give the per-oracle numbers for the
paper (Section~\ref{sec:ledger}).}

\Q{What happens to the unreproduced recovery failure?}{It stays in the
evidence folder; we report only what we can reproduce. If it reappears in a
campaign, triage will have a fresh chance with the full log.}

\subsection{LLM arm}

\Q{Why do you need an LLM at all?}{To keep the grammar current with code that
changes weekly. The model reads a commit and writes grammar that reaches the
changed code. Doing that by hand for 30 commits would take days.}

\Q{What if the LLM hallucinates?}{It cannot hurt the campaign: a fragment must
parse, be derivable, run on the server at $\geq$60\% validity, and (with a
coverage build) execute more of the commit's changed lines. A rewrite must
pass a value probe (NULL, type edges, data quantiles) and 150 live test cases
with no mismatch. Rejections go back to the model with the server's messages,
up to 3 rounds; everything is logged.}

\Q{Did the LLM ever propose something wrong that the gate caught?}{Yes, by
construction in our tests: \texttt{e + 1} was rejected with 12 mismatches;
\texttt{e IS TRUE} passed 109 sampled live checks but the value probe rejected
it at NULL. That is why the probe exists: sampling is not proof.}

\Q{Isn't there data leakage in RQ3?}{The miner excludes the 23 fixes we
revert, so the model never sees them in the prompt. Older public commits may
be in its training data; that is a threat to validity we state. The mining
window ends at the benchmark base commit.}

\Q{How much does it cost?}{About \$0.10 per round with a cached 9k-token
system prompt; \$5--10 for 30 commits with repair rounds (estimate; the
ledger will give the exact number). Manual mode costs nothing in API.}

\Q{What is diff coverage?}{Changed executable lines of the commit that the
generated SQL executed, divided by all changed executable lines, measured with
gcov on a coverage build at the same commit. It tells us whether a fragment
reaches the code it was written for.}

\Q{Why did the +LLM arm in the smoke run equal Grammar+Oracles?}{It ran
without fragments: the extension run is Phase 1 of the plan. It is a
correctness check of the pipeline: the arm must behave like its base.}

\subsection{Evaluation and statistics}

\Q{Why 5 trials? Klees says more.}{Budget: 7 arms $\times$ 5 trials $\times$ 1
hour is already 9 hours on 4 cores. We report medians, IQR, bootstrap CIs,
Mann--Whitney $p$ and $\hat A_{12}$, so the reader sees effect sizes even
where $p$ is not small. With 5 vs 5 the smallest exact two-sided $p$ is
0.008. If time allows we add trials.}

\Q{Why Mann--Whitney and not a t-test?}{Fuzzing metrics are skewed and
bounded (time to bug is censored at the budget); a rank test makes no
normality assumption. Arcuri \& Briand recommend it with $\hat A_{12}$ for
randomized algorithms.}

\Q{How do you handle bugs not found within the budget?}{Time-to-bug is
censored at 3600 s. We report, per bug, in how many trials each arm found it
and the median time among those trials, and test the per-trial count of bugs
found.}

\Q{Is the comparison fair to the random arm?}{It gets the same server,
oracles, schema, budget and seeds. It runs 10$\times$ more test cases (they are
short) and still reaches 17 features vs 101. Its weakness is validity: 26\%
of its statements run. That is the point of the comparison.}

\Q{Why is your benchmark valid ground truth?}{The bugs are real \pg bugs,
fixed by \pg developers with regression tests. We verify that each PoC shows a
symptom on its solo server and not on the fixed server, and the same symptom
with all 23 injected (no masking). We count only the 20 observable ones.}

\Q{What if two injected bugs interact?}{Verify shows no masking for the PoCs.
Attribution replays each finding on each solo server; a finding that needs
two bugs together is ``unattributed'' and not credited to any arm.}

\Q{How is coverage measured with 4 parallel trials?}{Each worker has its own
coverage server with its own \texttt{GCOV\_PREFIX}, so counters never merge;
lcov captures every 300 s.}

\Q{What are the threats to validity?}{One system, one commit, one schema;
benchmark of recent fixes; our own oracles and verdicts; the rewrite gate is
testing not proof; one LLM at one time; coverage is not detection; attribution
by replay (Section~\ref{sec:paper}).}

\subsection{Course connection and process}

\Q{Which course concepts does the project use?}{Grammar fuzzing and
black/grey-box testing (L14--16), coverage and mutation testing (L13: the
re-injected benchmark), soundness/precision (L2--L5: generator validity, oracle
precision), type systems (L29--30: typed grammar), interval and
constraint-based analysis (L19--20, L27--28: boundary and selective literals),
information flow (L21--22: RLS oracle), Hoare-style assertions (each rewrite is
an equivalence assertion), ASTs/CFGs (derivation trees, gcov). Table in
Section~\ref{sec:background}.}

\Q{Who did what?}{Venu: recent \pg features, campaigns, new grammar rules.
Bhuvan: first pass on every bug candidate (reproduce, shrink, decide).
Gowtham: experiment design, statistics, LLM arm, paper. All three know the
whole system; this guide is our shared script.}

\Q{How did you use AI tools?}{The LLM arm uses Claude as part of the method
(fragments, seeds, rewrites; every prompt and response stored). We also used
an AI coding assistant during development; all code is tested (138 unit tests)
and every reported number comes from a run in the repository. The paper's AI
statement discloses both.}

\Q{What would you do with two more months?}{More trials and longer campaigns;
mutation of the schema and data during a campaign; more oracles from the
plan-variant $\times$ rewrite product (e.g.\ rewrites under each variant);
LLM-proposed oracles beyond expression identities; other DBMSs with the same
grammar engine (only the grammar and executor are \pg-specific).}

\Q{What is the single biggest lesson?}{Oracle precision, not generator power,
limits a wrong-result fuzzer: a richer grammar finds the corners where two
``equivalent'' queries are not equivalent. Measuring that precision, and
fixing each false-positive class with a test, is what makes the reports
trustworthy.}
